\chapter{Módulo 4 --- Importar y ordenar datos (tidy data)}

\begin{itemize}
\tightlist
\item
  \textbf{Duración estimada:} 1.5 h teoría + 2.5 h práctica
\item
  \textbf{Objetivos de aprendizaje.} Al finalizar, la persona podrá:

  \begin{enumerate}
  \def\labelenumi{\arabic{enumi}.}
  \tightlist
  \item
    \textbf{Explicar} los tres principios de los datos ordenados
    (\emph{tidy data}).
  \item
    \textbf{Importar} archivos CSV y Excel con
    \passthrough{\lstinline!readr!} y \passthrough{\lstinline!readxl!},
    controlando tipos de columna.
  \item
    \textbf{Transformar} tablas entre formato ancho y largo con
    \passthrough{\lstinline!pivot\_longer()!} /
    \passthrough{\lstinline!pivot\_wider()!}.
  \item
    \textbf{Diferenciar} los problemas de calidad más comunes (tipos
    erróneos, \passthrough{\lstinline!NA!}, duplicados).
  \end{enumerate}
\end{itemize}

\section{Contenidos}

\begin{enumerate}
\def\labelenumi{\arabic{enumi}.}
\tightlist
\item
  El \emph{tidyverse} y los \emph{tibbles}.
\item
  Principios de datos ordenados: cada variable una columna, cada
  observación una fila, cada valor una celda.
\item
  \passthrough{\lstinline!readr::read\_csv()!}, especificación de tipos
  con \passthrough{\lstinline!col\_types!},
  \passthrough{\lstinline!problems()!}.
\item
  \passthrough{\lstinline!readxl::read\_excel()!}: hojas y rangos.
\item
  \passthrough{\lstinline!tidyr::pivot\_longer()!} y
  \passthrough{\lstinline!pivot\_wider()!}.
\item
  \passthrough{\lstinline!separate\_wider\_delim()!} y
  \passthrough{\lstinline!unite()!}.
\item
  Manejo de \passthrough{\lstinline!NA!}:
  \passthrough{\lstinline!drop\_na()!},
  \passthrough{\lstinline!replace\_na()!}; duplicados con
  \passthrough{\lstinline!distinct()!}.
\item
  Exportar: \passthrough{\lstinline!write\_csv()!}.
\end{enumerate}

\section{Desarrollo}

\subsection{Importar con readr}

\begin{lstlisting}[language=R]
library(tidyverse)

ventas <- read_csv("datos/ventas.csv", col_types = cols(
  fecha     = col_date(),
  tienda_id = col_character(),
  producto  = col_character(),
  precio    = col_double(),
  unidades  = col_integer()
))
ventas
glimpse(ventas)
tiendas <- read_csv("datos/tiendas.csv", show_col_types = FALSE)
\end{lstlisting}

\subsection{Excel}

\begin{lstlisting}[language=R]
library(readxl)
# excel_sheets("datos/archivo.xlsx")                 # lista las hojas
# read_excel("datos/archivo.xlsx", sheet = "Ventas", range = "A1:E500")
\end{lstlisting}

\subsection{De ancho a largo y de regreso}

Muchas hojas de cálculo guardan un mes por columna: eso \textbf{no} es
\emph{tidy}, porque ``mes'' es una variable.

\begin{lstlisting}[language=R]
reporte_ancho <- tribble(
  ~tienda, ~ene, ~feb, ~mar,
  "T01",    120,   98,  143,
  "T02",     87,  110,   95
)

reporte_largo <- reporte_ancho |>
  pivot_longer(cols = ene:mar, names_to = "mes", values_to = "unidades")
reporte_largo

reporte_largo |>
  pivot_wider(names_from = mes, values_from = unidades)
\end{lstlisting}

\begin{quote}
\passthrough{\lstinline!|>!} es el \emph{pipe} nativo de R (≥ 4.1):
\passthrough{\lstinline!x |> f(y)!} equivale a
\passthrough{\lstinline!f(x, y)!}. Léelo como ``y luego''.
\end{quote}

\subsection{Separar y unir columnas}

\begin{lstlisting}[language=R]
codigos <- tibble(sku = c("CEL-001-NEGRO", "TAB-014-GRIS"))
codigos |>
  separate_wider_delim(sku, delim = "-", names = c("categoria", "numero", "color"))
\end{lstlisting}

\subsection{Valores faltantes y duplicados}

\begin{lstlisting}[language=R]
ventas |> summarise(faltantes = sum(is.na(unidades)))
ventas_sin_na <- ventas |> drop_na(unidades)
ventas |> mutate(unidades = replace_na(unidades, 0L)) |> head()
ventas |> distinct(tienda_id, fecha, producto) |> nrow()   # verifica llave única
\end{lstlisting}

\section{Actividades y ejercicios}

\begin{enumerate}
\def\labelenumi{\arabic{enumi}.}
\tightlist
\item
  Importa \passthrough{\lstinline!ventas.csv!} con
  \passthrough{\lstinline!read\_csv()!} sin
  \passthrough{\lstinline!col\_types!} y compara el tipo que asigna a
  \passthrough{\lstinline!unidades!} y \passthrough{\lstinline!fecha!}
  con la versión especificada.
\item
  Convierte \passthrough{\lstinline!ventas!} a formato ancho: una fila
  por \passthrough{\lstinline!fecha!} y
  \passthrough{\lstinline!tienda\_id!}, una columna por producto con las
  unidades.
\item
  Regresa el resultado anterior a formato largo y comprueba que
  recuperas el mismo número de filas.
\item
  Cuenta los \passthrough{\lstinline!NA!} de
  \passthrough{\lstinline!unidades!} por producto.
\end{enumerate}

\subsection{Soluciones}

\begin{lstlisting}[language=R]
# 1
v2 <- read_csv("datos/ventas.csv", show_col_types = FALSE)
spec(v2)          # readr infiere fecha como date y unidades como double

# 2
ancho <- ventas |>
  select(fecha, tienda_id, producto, unidades) |>
  pivot_wider(names_from = producto, values_from = unidades)
ancho

# 3
largo <- ancho |>
  pivot_longer(-c(fecha, tienda_id), names_to = "producto", values_to = "unidades")
nrow(largo) == nrow(ventas)    # TRUE

# 4
ventas |>
  group_by(producto) |>
  summarise(n_na = sum(is.na(unidades)))
\end{lstlisting}

\section{Evaluación (formativa)}

\begin{itemize}
\tightlist
\item
  \textbf{Quiz:} Da un ejemplo de tabla que viole cada principio de
  \emph{tidy data}. ¿Qué función usarías para arreglarla?
\item
  \textbf{Práctica:} Recibe un Excel ``reporte mensual'' con meses en
  columnas y entrégalo en formato largo como CSV con
  \passthrough{\lstinline!write\_csv()!}.
\end{itemize}

\section{Referencias verificadas}

\begin{itemize}
\tightlist
\item
  \textbf{Artículo:} Wickham, H. (2014). Tidy data. \emph{Journal of
  Statistical Software, 59}(10), 1--23.
  \url{https://doi.org/10.18637/jss.v059.i10} {[}EN{]} {[}OA{]}
\item
  \textbf{Artículo:} Wickham, H., Averick, M., Bryan, J., Chang, W.,
  McGowan, L. D., François, R., \ldots{} Yutani, H. (2019). Welcome to
  the tidyverse. \emph{Journal of Open Source Software, 4}(43), 1686.
  \url{https://doi.org/10.21105/joss.01686} {[}EN{]} {[}OA, CC BY 4.0{]}
\item
  \textbf{Libro:} Wickham, H., Çetinkaya-Rundel, M., \& Grolemund, G.
  (2023). \emph{R for Data Science} (2.ª ed.). O'Reilly Media. ISBN
  978-1-4920-9740-2. Español: \url{https://davidrsch.github.io/r4ds-es/}
  {[}ES{]} {[}OA{]} --- caps. ``Ordenar datos'', ``Importación de
  datos'', ``Hojas de cálculo''.
\item
  \textbf{Video:} Gomila Salas, J. G. (Frogames). \emph{Curso completo
  de R para Data Science con Tidyverse -- R for Data Science en Español}
  {[}Video{]}. YouTube.
  \url{https://www.youtube.com/watch?v=9NJyhs5PlGQ} {[}ES{]}
  \emph{{[}nombre exacto del canal y fecha POR VERIFICAR{]}}
\end{itemize}
